\subsection{Efficiency and repeated independent matchings}
The disjoint-pair construction exchanges information for an exact conditional test. Randomized incomplete U-statistics offer related statistical/computational trade-offs \citep{chenkato2019}. This cost is explicit at a continuous independence model. For adjacent edges $e=\{v,a\}$ and $f=\{v,b\}$, write $h=H_e-H_f$. Then $\operatorname{Var}(h)=2$. After transforming the three margins to independent uniforms, its first-order projection is
\[
 g(U)=(2U_v-1)\{(2U_a-1)-(2U_b-1)\},\qquad
 \zeta_1=\operatorname{Var}\{g(U)\}=2/9.
\]
The paired estimator has variance $2/N\sim4/n$, while the full order-two U-statistic has asymptotic variance $4\zeta_1/n=8/(9n)$. Their variance ratio is $4.5$, corresponding to about $22\%$ relative efficiency. This calculation describes the continuous independence model; power also depends on the alternative, discordances and the simultaneous-testing family.

We therefore evaluate $R=8$ matchings, the first predetermined and the others independent uniform permutations chosen independently of the measurements. Independent auxiliary tie randomization is used in each matching. Let $p^{(r)}_{ef}$ be its unadjusted one-sided p-value and define
\begin{equation}\label{eq:merge}
 \widetilde p_{ef}=\min\left\{1,\frac{2}{R}\sum_{r=1}^R p^{(r)}_{ef}\right\}.
\end{equation}
For $R=1$ we retain the original p-value. The factor two is the dependence-safe arithmetic-mean merging rule of \citet{vovk2020}. Matchings share observations, so their p-values are dependent.

\begin{proposition}[Finite-sample validity of matching aggregation]\label{prop:merge}
Under the sampling assumptions of Theorem~\ref{thm:finite}, apply Holm to \eqref{eq:merge} over a fixed comparison family. Retain any subset of its rejections and then take transitive closure. The resulting strict orders, and every subsequent degree or rank projection, satisfy the same familywise guarantee $1-\alpha$.
\end{proposition}
\begin{proof}
Conditional on any independently chosen matching, each null p-value is superuniform by Theorem~\ref{thm:finite}. It remains superuniform after averaging over the matching randomization. The arithmetic-mean rule is valid under arbitrary dependence, so every merged null p-value is superuniform. Holm controls any false rejection. Deleting rejections cannot increase this event, and the transitive consequences of true strict orders are true.
\end{proof}
For reproducible closure, the implementation retains only merged rejections consistent with the average paired edge scores. This deletion preserves the guarantee. It does not introduce comparisons that failed Holm. Aggregating p-values reduces dependence on an individual matching, but its power must be measured: the factor two and p-value averaging can be conservative. For comparison, averaging $R$ independently randomized matching \emph{estimators} has variance $\operatorname{Var}(U)+(\operatorname{Var}(V)-\operatorname{Var}(U))/R$, by conditional expectation over matchings. At the independence model its asymptotic variance relative to the full U-statistic is $1+3.5/R$. This estimator calculation alone does not establish an equivalent efficiency for \eqref{eq:merge}.

\subsection{Exact evidence aggregation with deterministic tie handling}\label{sec:evalues}
P-value averaging can be conservative at a large comparison family. We therefore also aggregate nonnegative evidence with null expectation at most one, an e-value \citep{vovk2021}. For matching $r$, use the original tau-a kernels $K_{r,k,e}\in\{-1,0,1\}$, retaining zero at a tied coordinate, and set $Z_{r,k,ef}=(K_{r,k,e}-K_{r,k,f})/2\in[-1,1]$. For a predetermined fraction $0\le\lambda<1$, define
\[
 E_{r,\lambda,ef}=\prod_{k=1}^N(1+\lambda Z_{r,k,ef}).
\]
The evaluated multiset of seven fractions is
\[
 \Lambda_N=\{1/32,1/16,1/8,1/4,1/2,3/4,\min(1/2,N^{-1/2})\}.
\]
Its last element is determined by sample size and ensures consistency even for gaps too small for the fixed fractions. We compare $R=8$ and $R=64$ matchings, with the former nested in the latter, and put
\begin{equation}\label{eq:evalue}
 \overline E_{ef}=\frac1{7R}\sum_{r=1}^R\sum_{\lambda\in\Lambda_N}E_{r,\lambda,ef},\qquad
 p^E_{ef}=\min\{1,1/\overline E_{ef}\}.
\end{equation}
The averaging rule is classical e-value merging. Its use here combines exact Kendall contrast evidence across shared-observation matchings. The statistic handles ties directly, without auxiliary sign randomization. Matching randomization uses the same stream as the p-value comparator, with unused tie draws retained solely to align the matching sequences.

\begin{proposition}[Finite-sample evidence validity and fixed-family consistency]\label{prop:evalues}
Under IID rows, Holm applied to \eqref{eq:evalue} controls false strict orders at level $\alpha$, including discrete margins. Deleting rejections inconsistent with the full-sample Kendall scores and taking transitive closure preserves this control. For fixed $R$ and a fixed comparison family, each comparison with $\tau_e>\tau_f$ is rejected with probability tending to one as $n\to\infty$.
\end{proposition}
\begin{proof}
For a fixed independently selected matching, the contrasts $Z_k$ are IID with $\E Z_k=(\tau_e-\tau_f)/2$. Under the null this mean is nonpositive, so nonnegativity and independence give
\[
 \E E_{r,\lambda,ef}=\{1+\lambda\E Z_k\}^N\le1.
\]
Averaging over fractions and matchings preserves the expectation bound regardless of their dependence. Markov's inequality makes the reciprocal p-value superuniform. Holm, deletion and transitive closure give the common event used by all projections.

For a fixed positive gap $\Delta$, take the first matching and $\lambda_N=N^{-1/2}$ for $N\ge4$. Since $\log(1+u)\ge u-u^2$ for $|u|\le1/2$,
\[
 \log E_{1,\lambda_N,ef}\ge N^{-1/2}\sum_k Z_k-N^{-1}\sum_kZ_k^2
 \ge \sqrt N\{\Delta/2+o_{\Pp}(1)\}-1\ \longrightarrow\ \infty.
\]
Thus $\overline E_{ef}\ge E_{1,\lambda_N,ef}/(7R)$ diverges in probability, and $p^E_{ef}\to0$. The full-sample Kendall difference converges to $\Delta$, so the consistency filter retains the direction with probability tending to one. Holm's threshold is at least $\alpha/K>0$ for a fixed family, proving the assertion.
\end{proof}
The result establishes validity and consistency; relative power varies by design and is measured below. The simulations compare both evidence averages with the same eight-matching p-value pool and with full-U-statistic methods. For continuous margins and local comparisons, $H_{\{v,a\}}H_{\{v,b\}}=S_aS_b$: one $p\times p$ sign Gram matrix supplies all discordance counts in $O(Np^2+K)$ work per matching. Tied raw kernels use exact five-category contrast counts. Logarithmic accumulation prevents numerical overflow of products, while independent direct-kernel checks verify the counts and evidence calculations.

\subsection{A finite-sample comparator using all observation pairs}
Define $k_e(X_s,X_t)=\operatorname{sgn}(X_{si}-X_{ti})\operatorname{sgn}(X_{sj}-X_{tj})$ and
\[
 \widehat\tau_e={\binom n2}^{-1}\sum_{s<t} k_e(X_s,X_t),\qquad
 \widehat\Delta_{ef}=\widehat\tau_e-\widehat\tau_f.
\]
The difference kernel lies in $[-2,2]$, including ties. Hoeffding's order-two U-statistic inequality \citep{hoeffding1963} gives
\begin{equation}\label{eq:directhoeffding}
 \Pp\{\widehat\Delta_{ef}-\Delta_{ef}\ge t\}
 \le \exp(-Nt^2/8).
\end{equation}
Consequently $p^{U}_{ef}=\exp\{-N(\max\{\widehat\Delta_{ef},0\})^2/8\}$ is superuniform under $\Delta_{ef}\le0$. Holm and the same projections therefore give finite-sample simultaneous inference. To see \eqref{eq:directhoeffding} directly, average the disjoint-pair mean over a uniform random matching conditional on the sample. This conditional average is $\widehat\Delta_{ef}$. Jensen's inequality bounds its exponential moment by that of a mean of $N$ independent kernels of range four; optimizing the usual bounded-variable exponential bound gives \eqref{eq:directhoeffding}. Thus dependence between the $\binom n2$ kernels is respected.

This comparator isolates direct testing from rectangularization while retaining a finite-sample guarantee. Its bounded-range calibration can remain conservative even though its estimator uses all pairs.

\subsection{Full U-statistics with Gaussian multiplier calibration}\label{sec:multiplier}
For a computationally stronger comparator, estimate the first-order influence by
\[
 \widehat\psi_{te}=2\left\{\frac1{n-1}\sum_{s\ne t} k_e(X_s,X_t)-\widehat\tau_e\right\}.
\]
Given independent standard normal multipliers $\xi_t^{(b)}$, put $Z_e^{*(b)}=n^{-1}\sum_t\xi_t^{(b)}\widehat\psi_{te}$ and
\[
 s_{ef}^2=n^{-2}\sum_t(\widehat\psi_{te}-\widehat\psi_{tf})^2,\qquad
 Q_b=\max_{\{e,f\}\in\mathcal C}\frac{|Z_e^{*(b)}-Z_f^{*(b)}|}{s_{ef}}.
\]
The empirical $1-\alpha$ quantile $q$ gives intervals $\widehat\Delta_{ef}\pm qs_{ef}$. Intervals excluding zero supply strict orders. This is an asymptotic IID procedure for nondegenerate projections, grounded in Gaussian and multiplier approximation for high-dimensional U-statistics \citep{chen2018,cck2013}. Theorem~\ref{thm:finite} applies to the exact methods above; multiplier calibration has its own approximation assumptions. Numerically degenerate contrasts are omitted, and their counts are reported.

The matched multiplier rectangular comparator uses the same scores, influences and multiplier draws, calibrates the maximum standardized error across individual edges, and converts its edge intervals into strict orders over the same family $\mathcal C$. Both comparators then use identical tree projections. This comparison measures the effect of calibrating differences rather than separate edge weights.

Tie-safe forward and backward Fenwick sweeps compute every rowwise kernel mean in $O(p^2n\log n)$ time. Contrast maxima are streamed in C++ without allocating a $B\times|\mathcal C|$ matrix. These implementations retain the full tau-a statistic and its shared-edge covariance.

\subsection{Local and complete families}
Every fixed comparison family can be used in the exact constructions. We compare all pairs sharing a vertex with all pairs of distinct edges. Their ordered hypothesis counts on a complete graph are $K_{\rm local}=p(p-1)(p-2)$ and $K_{\rm all}=m(m-1)$, with ratio $(p+1)/4$ for $p\ge3$. At $p=52$ these are $132{,}600$ and $1{,}756{,}950$.

Complete comparisons can add nonlocal information relevant to cycle-path constraints, while also increasing the multiplicity penalty and computation. Transitive closure of local orders recovers some nonlocal comparisons, but need not recover all of them. The family ablation measures the resulting trade-off at $p=12$; both families are also evaluated on every cross-sectional application. Local testing remains an explicit computational choice for the larger simulations and the Polish sensitivity study.
